% Pursuing Stacks (cotes 134-2 to 134-8): archive cote 134-4, site page 19 (PDF page 20), typescript page 107.
% typescript. Draft compared: ps3.tex:928-990.
% AI-transcribed by Claude Opus 5.5 (claude-opus-5-5, adaptive thinking, effort high), Message Batches API, batch ps1 (msgbatch_01Rv5hCaNTtUt4JHx1yqoHPL), 2026-11-12; not yet checked by eye.
\scanpage{134-4/19}
\mspage{107}
being aspheric ),  and in the case when M is one among the usual elementa-%
ry modelizers.  The question of devising notions which still make sense
when there is no homotopy interval in M is'nt perhaps so silly after all,
if we remember that the category of semi-simplicial face complexes (name-%
ly without degeneracies) is indeed a modelizer,  but it has'nt got a homo-%
topy interval.  However for the time being I feel it \unsure{is'nt|isn't} too urgent to
\margin{$\lfloor$} get any more into this.

\margin{$\lceil$} 49.   I would like now to elaborate a little on the notion of a contractib
le element  $X$ in $M$,  which \inserted[typed]{(I recall)} means an object ,  admitting a section $e_o$:$e\to X$,
such that the constant map
\[ c = p_X e_o \ :\ X \to e \to X \]
is  homotopic to $\mathrm{id}_X$  ,  i.e. there exists a homotopy interval I such that
$\mathrm{id}_X \underset{\underline{I}}{\sim} c$  ($\underline{I}$-homotopy in the wide sense).

If $X$ is contractible,  then the constant map $c$: $X\to X$ is a weak equiva-%
lence (as it is homotopic to $\mathrm{id}_X$ which is )  and hence by the satturation
condition c') (in fact the mild saturation condition c) suffices) it fol-%
lows that $p_X$: $X\to e$ is a weak equivalence.  In fact, \inserted{one would expect it} \deleted{to be} even universal
-ly so,  i.e. \xout{xxxxx}, \deleted{an} \inserted{that px is} \uline{aspheric} map,  as a plausible generalization of the
homotopy lemma c)  above (which we have used already a number of times as
our main asphericity criterion in elementary modelizers $A\hat{}$...).  The natu-%
ral idea to prove asphericity of $p_X$ ,  namely that for any base change
$S\to e$ ,  the projection \inserted[typed]{$X_S=$} $X$x$S\to S$ is a weak equivalence,  is to apply
the precedent criterion  to $X_S$ , viewed as an element of $M_{/S}$ .  As the
base change functor commutes with products,  it is clear indeed that
for two maps $f,g$:  $X \rightrightarrows Y$  in $M$,  $\underline{I}$-homotopy of $f$ and $g$ will imply
$\underline{I}_S$-homotopy of $f_S$ and $g_S$ ,  with evident definition of base change for an
``interval''.  On the other hand,  asphericity of $\underline{I}$ over $e$ implies tautolo-%
gically asphericity of $\underline{I}_S$ over $S$,  the final object of $M_{/S}$ ,  which is all
we need to care about to get the result that $X_S\to S$ is a weak equivalence.
(The condition b) on homotopy intervals is definitely misleading here,  as
it would induce us to put the extra condition that the base change functor
is compatible with initial elements,  which is true indeed if $\emptyset_M$ is strict\hand{,}
but here a wholly extraneous condition ...)

Thus $X$ \uline{contractible implies} $X$ \uline{aspheric over} $e$ .  It is well known
that the converse is'nt true,  already for  $M$ = (Ss sets) , or
$M$ = (Sp\over{\illegible{1}}{a}ces)  with the usual notions of weak equivalences,  t\over{\unsure{q}}{a}king aspheric
complexes which are not Khan complexes,  or``aspheric''spaces (in the sense
of singular cohomology )  which are not CW-spaces.  The same examples
show, too,  that even a good honest homot\over{a}{o}py interval need not be contrac-%

\margin{19}
